% =====================================================================
%  Section 8.1--8.3: Cole--Hopf, exact solution, vanishing viscosity.
%  Section 8.4 (Godunov scheme) is in ext_godunov.tex.
% =====================================================================

\section{Extensions}
\label{sec:extensions}

This section tests the entropy condition against the two standard selection
mechanisms: vanishing viscosity (Subsections~\ref{subsec:cole-hopf}--%
\ref{subsec:vanishing-viscosity}) and monotone numerical approximation
(Subsection~\ref{subsec:godunov}).

\subsection{The Cole--Hopf transformation}
\label{subsec:cole-hopf}

\begin{theorem}[Cole--Hopf transformation \cite{cole1951,hopf1950}]
\label{thm:cole-hopf}
Let $u\in C^{2,1}(\mathbb{R}\times(0,T))$ solve $u_t+u\,u_x=\nu\,u_{xx}$, with $u$ and
$u_x$ continuous on $\mathbb{R}\times[0,T)$ (so that $c$ below is defined), and set
\begin{equation}
  \label{eq:cole-hopf-def}
  \hat\psi(x,t)=-\frac{1}{2\nu}\int_0^x u(y,t)\,dy,
  \qquad
  c(t)=\int_0^t\Big(\frac{u(0,\tau)^2}{4\nu}-\frac{u_x(0,\tau)}{2}\Big)d\tau,
  \qquad
  \theta(x,t)=\exp\big(\hat\psi(x,t)+c(t)\big)>0 .
\end{equation}
Then $\theta$ solves the heat equation $\theta_t=\nu\,\theta_{xx}$, and
$u=-2\nu\,\theta_x/\theta$.
\end{theorem}

\begin{proof}
Since $c$ depends on $t$ only, $\theta_x/\theta=\hat\psi_x=-u/(2\nu)$, which
gives $u=-2\nu\theta_x/\theta$. For the heat equation, let
$w=\hat\psi_t-\nu\hat\psi_{xx}-\nu\hat\psi_x^2$. Using
$\hat\psi_x=-u/(2\nu)$, $\hat\psi_{xx}=-u_x/(2\nu)$,
$\hat\psi_{xt}=-u_t/(2\nu)$ and $\hat\psi_{xxx}=-u_{xx}/(2\nu)$,
\begin{equation}
  w_x=\hat\psi_{xt}-\nu\hat\psi_{xxx}-2\nu\hat\psi_x\hat\psi_{xx}
  =-\frac{1}{2\nu}\big(u_t+u\,u_x-\nu\,u_{xx}\big)=0 .
\end{equation}
Hence $w$ depends on $t$ only and may be evaluated at $x=0$, where
$\hat\psi(0,t)=0$ gives $\hat\psi_t(0,t)=0$:
\begin{equation}
  w(t)=w(0,t)=\frac{u_x(0,t)}{2}-\frac{u(0,t)^2}{4\nu}=:C(t).
\end{equation}
With $\psi=\hat\psi+c$ we have $c'=-C$, so $\psi_t-\nu\psi_{xx}-\nu\psi_x^2=
C+c'=0$. Since $\theta=e^{\psi}$ satisfies $\theta_t=\theta\psi_t$ and
$\theta_{xx}=\theta(\psi_{xx}+\psi_x^2)$, this is $\theta_t=\nu\theta_{xx}$.
\end{proof}

\begin{remark}
The factor $e^{c(t)}$ is independent of $x$, so it affects neither
$u=-2\nu\theta_x/\theta$ nor the initial data ($c(0)=0$). It cannot be
omitted: $e^{\hat\psi}$ satisfies $\partial_t\theta=\nu\theta_{xx}+C(t)\theta$,
which is the heat equation only if $C\equiv0$.
\end{remark}

\subsection{Exact solution via the heat kernel}
\label{subsec:cole-hopf-solution}

By Theorem~\ref{thm:cole-hopf}, a classical solution of \eqref{eq:viscous-burgers}
that is bounded on $\mathbb{R}\times[0,T']$ for each $T'<T$ (this is what makes
$\theta$ obey the growth bound below) corresponds to a solution of the heat
equation with initial data
\begin{equation}
  \label{eq:theta0}
  \theta_0(x) = \exp\left(-\frac{1}{2\nu}\int_0^x u_0(y)\, dy\right)
  \le e^{\|u_0\|_\infty|x|/(2\nu)} .
\end{equation}
The Cauchy problem for the heat equation has a unique solution in the class of
at most $e^{a|x|^2}$ growth \cite[\S2.3.3]{evans2010}, which contains this
data; it is given by
\begin{equation}
  \label{eq:theta-heat-kernel}
  \theta(x,t) = \int_{\mathbb{R}} \frac{1}{\sqrt{4\pi\nu t}}\,
    \exp\left(-\frac{(x-y)^2}{4\nu t}\right) \theta_0(y)\, dy.
\end{equation}
Substituting \eqref{eq:theta0}--\eqref{eq:theta-heat-kernel} into
$u=-2\nu\,\theta_x/\theta$ and differentiating under the integral gives the
\emph{Cole--Hopf formula}
\begin{equation}
  \label{eq:cole-hopf-formula}
  u(x,t) = \frac{\displaystyle\int_{\mathbb{R}} \frac{x-y}{t}\, \exp\left(-\frac{G(x,y,t)}{2\nu}\right) dy}
                {\displaystyle\int_{\mathbb{R}} \exp\left(-\frac{G(x,y,t)}{2\nu}\right) dy},
  \qquad
  G(x,y,t) = \int_0^y u_0(r)\, dr + \frac{(x-y)^2}{2t}.
\end{equation}

\begin{corollary}[Global smoothness; Hopf \cite{hopf1950}]
\label{cor:global-smoothness}
For every bounded measurable $u_0$ and every $\nu>0$, formula
\eqref{eq:cole-hopf-formula} defines a $C^\infty$ solution $u^{(\nu)}$ of
\eqref{eq:viscous-burgers} on $\mathbb{R}\times(0,\infty)$. Moreover
$\|u^{(\nu)}\|_\infty\le\|u_0\|_\infty$, $u^{(\nu)}(\cdot,t)\to u_0$ in
$L^1_{\mathrm{loc}}$ and at every continuity point of $u_0$ as $t\to0^+$, and
$u^{(\nu)}$ is the unique bounded solution with these properties.
\end{corollary}

\noindent For $t>0$, $\theta$ in \eqref{eq:theta-heat-kernel} is smooth and
positive, and for any such solution $\theta$ of the heat equation
$u=-2\nu\theta_x/\theta$ solves \eqref{eq:viscous-burgers}: with
$\psi=\ln\theta$, $\psi_t-\nu\psi_{xx}-\nu\psi_x^2=0$, and differentiating in
$x$ gives $(-u/2\nu)_t-\nu(-u/2\nu)_{xx}-\nu\big((u/2\nu)^2\big)_x=0$, i.e.\
$u_t+u\,u_x=\nu u_{xx}$. This proves the first assertion for merely measurable
$u_0$ (Theorem~\ref{thm:cole-hopf} is not needed in this direction), while the
bound, the sense in which the data are attained, and uniqueness are taken from
Hopf \cite{hopf1950}.

\subsection{The vanishing viscosity limit}
\label{subsec:vanishing-viscosity}

\begin{theorem}[Vanishing viscosity limit \cite{hopf1950}]
\label{thm:vanishing-viscosity}
Let $u_0\in L^\infty(\mathbb{R})$, let $u^{(\nu)}$ be the solution of
\eqref{eq:viscous-burgers} with data $u_0$ given by
\eqref{eq:cole-hopf-formula}, and let $u$ be the entropy solution of
\eqref{eq:inviscid-burgers} with the same data
(Definition~\ref{def:entropy-solution}; for Riemann data,
Theorem~\ref{thm:riemann-solution}). Then $u^{(\nu)} \to u$ pointwise almost
everywhere and in $L^1_{\mathrm{loc}}$ as $\nu \to 0^+$.
\end{theorem}

Hopf's proof \cite{hopf1950} applies Laplace's method to
\eqref{eq:cole-hopf-formula} and shows that $u^{(\nu)}$ converges to the
Lax--Oleinik solution \eqref{eq:lax-oleinik}. Because his work predates Oleinik's
condition, the identification of this limit with $u$ combines the theorem with
Remark~\ref{rem:lax-oleinik} below, which shows that the Lax--Oleinik solution
satisfies \eqref{eq:oleinik}, and with Theorem~\ref{thm:oleinik-uniqueness}.

\begin{remark}[The Lax--Oleinik formula]
\label{rem:lax-oleinik}
As $\nu\to0$ the weight $\exp(-G/2\nu)$ in \eqref{eq:cole-hopf-formula}
concentrates at the minimisers of $y\mapsto G(x,y,t)$. Where the minimiser
$y^*(x,t)$ is unique, the limit is
\begin{equation}
  \label{eq:lax-oleinik}
  u(x,t)=\frac{x-y^*(x,t)}{t},\qquad
  y^*(x,t)=\operatorname*{arg\,min}_y\Big\{V(y)+\frac{(x-y)^2}{2t}\Big\},
  \quad V(y)=\int_0^y u_0 ,
\end{equation}
the \emph{Lax--Oleinik formula} \cite{lax1957,oleinik1957},
\cite[\S3.4.2]{evans2010}; a jump occurs where two minimisers tie. Since
$x\mapsto y^*(x,t)$ is nondecreasing \cite[\S3.4.2]{evans2010}, for $z>0$
\begin{equation*}
  u(x+z,t)-u(x,t)=\frac{z-\big(y^*(x+z,t)-y^*(x,t)\big)}{t}\ \le\ \frac zt,
\end{equation*}
which is \eqref{eq:oleinik}. As \eqref{eq:lax-oleinik} also defines a weak
solution \cite[\S3.4.2]{evans2010}, the Lax--Oleinik solution is an entropy
solution in the sense of Definition~\ref{def:entropy-solution}. The same bound
shows $u(x+,t)-u(x-,t)\le0$ wherever one-sided limits exist: every jump is
downward. The minimiser is unique for almost every $(x,t)$
\cite[\S3.4.2]{evans2010}, so \eqref{eq:lax-oleinik} defines $u$ a.e.
\end{remark}

\begin{remark}[The Riemann problem from the Lax--Oleinik formula]
\label{rem:riemann-lax-oleinik}
For the data \eqref{eq:general-riemann-data}, $V(y)=u_\ell y$ for $y<0$ and
$V(y)=u_r y$ for $y>0$, so $G(y)=V(y)+(x-y)^2/2t$ is convex on each half-line, with
critical points $y_\ell=x-u_\ell t$ (admissible iff $x\le u_\ell t$) and
$y_r=x-u_rt$ (iff $x\ge u_rt$), and a kink at $y=0$ with one-sided derivatives
$u_\ell-x/t$ and $u_r-x/t$. If $u_\ell<u_r$ and $u_\ell t<x<u_rt$, only the kink
is a local minimiser and $u=x/t$; outside this range the single admissible
critical point gives $u_\ell$ or $u_r$, in agreement with
\eqref{eq:riemann-rarefaction}.
If $u_\ell>u_r$ and $u_rt\le x\le u_\ell t$, both $y_\ell$ and $y_r$ are admissible
and the kink is a local maximum, with
$G(y_\ell)=u_\ell x-u_\ell^2t/2$ and $G(y_r)=u_rx-u_r^2t/2$, so
$G(y_\ell)-G(y_r)=(u_\ell-u_r)\big(x-\tfrac{u_\ell+u_r}2t\big)$. The tie, hence the
jump, is at $x=\tfrac{u_\ell+u_r}2t$, recovering \eqref{eq:riemann-shock} and the
Rankine--Hugoniot speed.
\end{remark}

\subsubsection{Viscous shock profiles}
\label{subsubsec:viscous-profiles}

\begin{proposition}[Viscous shock profiles]
\label{prop:viscous-profile}
Let $u_\ell\neq u_r$ and $\nu>0$. Equation \eqref{eq:viscous-burgers} has a
travelling-wave solution $u(x,t)=\phi(x-st)$, with $\phi\in C^2$ bounded,
$\phi(\eta)\to u_\ell$ as $\eta\to-\infty$, $\phi(\eta)\to u_r$ as $\eta\to+\infty$
and $\phi'\to0$ as $\eta\to\pm\infty$, if and only if $u_\ell>u_r$. In that case
$s=(u_\ell+u_r)/2$ and, with $\Delta=u_\ell-u_r$, $\phi$ is unique up to
translation $\eta\mapsto\eta+\eta_0$; the translate with $\phi(0)=s$ is
\begin{equation}
  \label{eq:viscous-profile}
  \phi(\eta)=s-\frac{\Delta}{2}\tanh\!\Big(\frac{\Delta\,\eta}{4\nu}\Big).
\end{equation}
Moreover $\|\phi-(s-\tfrac\Delta2\operatorname{sgn}\eta)\|_{L^1(\mathbb{R})}
=4\ln2\;\nu$, independent of $s$ and $\Delta$.
\end{proposition}

\begin{proof}
Substituting $u=\phi(x-st)$ gives $-s\phi'+\phi\phi'=\nu\phi''$, and
integrating once from $-\infty$ (where $\phi'\to0$),
\begin{equation*}
  \nu\phi'=\tfrac12\phi^2-s\phi-K,\qquad K=\tfrac12u_\ell^2-su_\ell .
\end{equation*}
Letting $\eta\to+\infty$ shows that the right-hand side must also vanish at
$u_r$, so $u_\ell,u_r$ are the two roots of the quadratic: $s=(u_\ell+u_r)/2$
(the Rankine--Hugoniot speed) and
\begin{equation}
  \label{eq:profile-ode}
  \nu\phi'=\tfrac12(\phi-u_\ell)(\phi-u_r).
\end{equation}
This autonomous first-order ODE cannot be crossed at its equilibria
$u_\ell,u_r$, so a connecting $\phi$ lies strictly between them. On
$(\min(u_\ell,u_r),\max(u_\ell,u_r))$ the right side of \eqref{eq:profile-ode}
is negative, so $\phi$ is decreasing. If $u_\ell<u_r$ it would have to
increase from $u_\ell$ to $u_r$, a contradiction; if $u_\ell>u_r$,
\eqref{eq:profile-ode} is solved by \eqref{eq:viscous-profile}; any other solution is a translate, by
uniqueness for the ODE. For the $L^1$
distance, by oddness and $\int_0^\infty(1-\tanh\zeta)\,d\zeta=\ln2$,
\begin{equation*}
  \int_{\mathbb{R}}\tfrac\Delta2\Big|\tanh\tfrac{\Delta\eta}{4\nu}-\operatorname{sgn}\eta\Big|\,d\eta
  =\Delta\int_0^\infty\Big(1-\tanh\tfrac{\Delta\eta}{4\nu}\Big)d\eta
  =\Delta\cdot\frac{4\nu}{\Delta}\ln2=4\ln2\,\nu .
\end{equation*}
\end{proof}

Hence a jump is the limit of viscous travelling waves, of thickness
$O(\nu/\Delta)$, exactly when $u_\ell>u_r$. This motivates the Lax condition
\eqref{eq:lax-burgers} from vanishing viscosity. That $u_A$ of
Section~\ref{sec:nonuniqueness}, with its upward jump, is not the
vanishing-viscosity limit is the content of Theorem~\ref{thm:vanishing-viscosity}
together with Proposition~\ref{prop:entropy-selects-B}.

For Riemann data the Cole--Hopf solution is available in closed form since $G$ is
piecewise quadratic (\texttt{code/cole\_hopf\_riemann.py}). For the shock $(1,0)$ its
$L^1(-2,2)$ distance to the entropy solution at $t=1$ equals $2.773\,\nu$ for
$\nu\le10^{-3}$, the layer value $4\ln2\,\nu$ of
Proposition~\ref{prop:viscous-profile}; for the rarefaction $(0,1)$ the error
divided by $\nu$ grows by about $4.4$ per decade of $\nu$ (observed, not derived).